\clearpage
\fancyhead[L]{\small JUNO.}
\hypertarget{native-255-lower}{}
\begin{center}
\textbf{ÜBER EINEN NEUEN, VON INSPEKTOR HARDING IN LILIENTHAL ENTDECKTEN HÖCHST MERKWÜRDIGEN WANDELSTERN.}
\end{center}

\begin{center}
Monatliche Correspondenz. Band X. S. 371–385 Oktober 1804.
\end{center}

\begin{center}
— — —
\end{center}

Kaum hatte ich dem unermüdlichen und unvergleichlichen Dr. \textsc{Gauss} meine ersten drei Beobachtungen dieses Gestirns vom 13., 14. und 15. September mitgeteilt, da erhielt ich am 23. September schon mit umgehender Post die folgende Antwort, die alle unsere Leser ebenso in Erstaunen versetzen wird wie mich. „Was werden Sie sagen“, schreibt dieser tiefsinnige Geometer, „dass ich es gewagt habe, meine eigenen Beobachtungen mit den drei Beobachtungen, die Sie mir freundlicherweise mitgeteilt haben, und ein paar früheren von Dr. \textsc{Olbers} zusammenzunehmen, die alle zusammen nur 14 Tage und einen heliozentrischen Bogen von vier Grad umfassen, dass ich es gewagt habe, auf diese wacklige Grundlage schon einen vorläufigen Versuch und elliptische Elemente einer Bahn ohne jede hypothetische Voraussetzung zu stützen? Das Ergebnis kann nur sehr \textit{prekär} sein. Doch ich bin geneigt zu hoffen, dass es nicht mehr \textit{enorm} oder \textit{völlig} von der Wahrheit abweichen kann, sondern wenigstens schon eine grobe Vorstellung von den Dimensionen der Bahn gibt. Mit noch größerer Zuversicht hoffe ich, dass es ausreichen wird, um danach den Planeten vielleicht einen Monat lang, vielleicht noch länger, aufzufinden. Und mit Gewissheit kann ich behaupten, dass dadurch alle bisherigen Beobachtungen gut dargestellt werden. Hier einstweilen das Ergebnis. Demnächst folgt der Vergleich mit den Beobachtungen, bei dem ich zugleich entscheiden werde, ob ich es des Titels \textit{Elemente I des Planeten von \textsc{Harding}} für würdig erklären kann.

\begin{center}
\begingroup\fontsize{9}{11}\selectfont\setlength{\tabcolsep}{3pt}
\begin{tabular}{lr}
Epoche, Seeberger Meridian, 1804, 5. Sept.  & $24^\circ\ 53'\ 44''$ \\
Sonnenferne . . . . . . . . . . & $244\ 51\ 36$ \\
aufsteigender Knoten . . . . . . & $171\ 48\ 24$ \\
Exzentrizität . . . . . . . . . & $0.313757$ \\
große Halbachse . . . . . . . . & $2.88208$ \\
tägliche Bewegung . . . . . . . & $725''18$ \\
Neigung der Bahn . . . . . . . & $15^\circ\ 12'\ 39''$ \\
Bewegung . . . . . . . . . . & rechtläufig \\
\end{tabular}\endgroup
\end{center}

„Was sagen Sie zu dieser sonderbaren Bahn, der großen Exzentrizität, der großen Annäherung an Ceres und Pallas hinsichtlich der Achse und der mittleren Bewegung, \textit{die durch eine kleine Änderung der Beobachtung sehr leicht zu völliger Gleichheit werden kann}? Ich will aber meinem Grundsatz treu bleiben, den Rechnungen schlechterdings nichts Hypothetisches beizumischen und künftigen Erfahrun%
\SourceLeaf{256}%
gen nicht vorzugreifen. In sehr kurzer Zeit werden wir schon viel weiter sein. Dass die Bahn himmelweit von einer Parabel verschieden ist und \textsc{Harding}s Stern die Bezeichnung Planet verdient, daran lässt sich jetzt schon kaum mehr zweifeln. Es wäre daher zu wünschen, dass er bald einen Namen erhält. Natürlich muss das Recht, ihn zu taufen, dem Entdecker allein vorbehalten bleiben usw.

Hier sind die Beobachtungen, die Dr. \textsc{Gauss} in Braunschweig angestellt hat.

\begin{center}
\begingroup\fontsize{9}{11}\selectfont\setlength{\tabcolsep}{3pt}\renewcommand{\arraystretch}{1.12}
\begin{tabular}{r|r|r|r}
1804 & \shortstack{Mittlere Zeit in\\Braunschweig} & \shortstack{Scheinbare gerade\\Aufsteigung} & \shortstack{Scheinbare südliche\\Deklination} \\ \hline
12. Sept. & $10^{u}\,35^{m}\,2^{s}$ & $0^\circ\ 54'\ 26''$ & $1^\circ\ 38'\ 15''$ \\
13 & $9\ 41\ 32$ & $0\ 45\ 24$ & $1\ 50\ 59$ \\
14 & $11\ 38\ 49$ & $0\ 35\ 37$ & $2\ 5\ 1$ \\
15 & $10\ 16\ 17$ & $0\ 26\ 53$ & $2\ 17\ 42$ \\
16 & $10\ 37\ 4$ & $0\ 17\ 17$ & $2\ 31\ 20$ \\
17 & $11\ 28\ 59$ & $0\ 7\ 23$ & $2\ 44\ 7$ \\
\end{tabular}\endgroup
\end{center}

\begin{center}
\begingroup\fontsize{9}{11}\selectfont\setlength{\tabcolsep}{3pt}\renewcommand{\arraystretch}{1.12}
\begin{tabular}{r|r|r|r}
1804 & \shortstack{Mittlere Zeit in\\Braunschweig} & \shortstack{Scheinbare gerade\\Aufsteigung} & \shortstack{Scheinbare südliche\\Deklination} \\ \hline
18 & $11^{u}\,22^{m}\,16^{s}$ & $359^\circ\ 57'\ 26''$ & $2^\circ\ 58'\ 3''$ \\
21 & $10\ 24\ 52$ & $359\ 28\ 6$ & $3\ 37\ 59$ \\
24 & $10\ 1\ 2$ & $358\ 57\ 53$ & $4\ 18\ 6$ \\
25 & $8\ 44\ 25$ & $358\ 48\ 12$ & $4\ 30\ 44$ \\
27 & $10\ 20\ 29$ & $358\ 27\ 20$ & $4\ 57\ 47$ \\
28 & $8\ 29\ 4$ & $358\ 18\ 20$ & $5\ 10\ 22$ \\
\end{tabular}\endgroup
\end{center}

Er schreibt dazu: „Die ersten drei Beobachtungen wurden mit einem schlechten und besonders schlecht montierten Achromaten gemacht und verdienen daher wenig Vertrauen. Die folgenden dagegen wurden mit einem sehr guten Spiegelteleskop gemacht. Sie werden daher besser sein, wenigstens so gut, wie es die Ringmikrometermethode und meine Sehschärfe zulassen.“

Wenige Tage nach diesem Schreiben, schon am 30. September, erhielt ich die nähere Bestätigung der Elemente dieser Planetenbahn. „Ich schicke Ihnen hier“, schreibt Dr. \textsc{Gauss}, „neue und verbesserte Elemente. Ich gebe mich der Hoffnung hin, dass sie schon eine annähernde Bestimmung der wahren Elemente liefern können, und wage es daher, sie als die \textit{I. Elemente} anzukündigen.

\begin{center}
\begingroup\fontsize{9}{11}\selectfont\setlength{\tabcolsep}{3pt}
\begin{tabular}{lr}
Epoche, Seeberger Meridian, 1804, 5. Sept. . . . . . . . & $20^\circ\ 38'\ 56''$ \\
Sonnenferne . . . . . . . . . . . . . . . . & $239\ 14\ 2$ \\
aufsteigender Knoten . . . . . . . . . . . . . & $171\ 15\ 35$ \\
Exzentrizität . . . . . . . . . . . . . . . . & $0.287359$ \\
große Halbachse (Logarithmus) . . . . . . . . . . & $0.438682$ \\
tägliche Bewegung . . . . . . . . . . . . . . & $779''80$ \\
Neigung der Bahn . . . . . . . . . . . . . . & $13^\circ\ 34'\ 59''$ \\
Bewegung . . . . . . . . . . . . . . . . . & rechtläufig. \\
\end{tabular}\endgroup
\end{center}

Mit diesen neuen Elementen verglich Dr. \textsc{Gauss} nun sämtliche Beobachtungen aus Seeberg, Bremen und Braunschweig. Die Differenzen sind nicht größer als die bei Ringmikrometern möglichen Fehler, die zum Teil auch auf die Vergleichssterne zurückgehen könnten.

\begin{center}
\textit{Seeberger Beobachtungen:}
\end{center}

\begin{center}
\begingroup\fontsize{8}{10}\selectfont\setlength{\tabcolsep}{3pt}\renewcommand{\arraystretch}{1.12}
\begin{tabular}{r|r|r|r|r|r}
1804 & \shortstack{Mittlere Zeit\\auf dem Seeberg} & \shortstack{Berechnete gerade\\Aufsteigung} & \shortstack{Differenz\\R—B} & \shortstack{Berechnete\\Deklination} & \shortstack{Differenz\\R—B} \\ \hline
13. Sept. & $12^{u}\,31^{m}\,59^{s}$ & $0^\circ\ 44'\ 57''5$ & $+1''0$ & $1^\circ\ 52'\ 34''7$ & $-2''3$ \\
14 & $12\ 27\ 27$ & $0\ 35\ 43.8$ & $+2.8$ & $2\ 5\ 35.3$ & $-0.3$ \\
15 & $12\ 22\ 53$ & $0\ 26\ 20.4$ & $+0.4$ & $2\ 18\ 40.5$ & $+2.0$ \\
17 & $12\ 13\ 45$ & $0\ 7\ 7.8$ & $-0.3$ & $2\ 45\ 8.7$ & $-4.4$ \\
18 & $12\ 9\ 10$ & $359\ 57\ 20.4$ & $0$ & $2\ 58\ 28.3$ & $-2.6$ \\
20 & $ .\quad .\quad . $ & $359\ 37\ 28.7$ & $ .\ .\ . $ & $3\ 25\ 14.9$ & $+11.4$ \\
23 & $11\ 46\ 11$ & $359\ 7\ 11.9$ & $-3.0$ & $4\ 5\ 32.0$ & $-2.5$ \\
\end{tabular}\endgroup
\end{center}
\Needspace{23\baselineskip}
\SourceLeaf{257}%
\begin{center}
\textit{Bremer Beobachtungen:}
\end{center}

\begin{center}
\begingroup\fontsize{8}{10}\selectfont\setlength{\tabcolsep}{3pt}\renewcommand{\arraystretch}{1.12}
\begin{tabular}{r|r|r|r|r|r}
1804 & \shortstack{Mittlere Zeit\\in Bremen} & \shortstack{Berechnete gerade\\Aufsteigung} & \shortstack{Differenz\\R—B} & \shortstack{Berechnete\\Deklination} & \shortstack{Differenz\\R—B} \\ \hline
7. Sept. & $10^{u}\,37^{m}\,21^{s}$ & $1^\circ\ 36'\ 57''$ & $+5''$ & $0^\circ\ 36'\ 6''$ & $-3''$ \\
8 & $8\ 11\ 20$ & $1\ 29\ 42$ & $+16$ & $0\ 47\ 4$ & $-15$ \\
9 & $10\ 48\ 50$ & $1\ 20\ 36$ & $+10$ & $1\ 0\ 47$ & $-3$ \\
10 & $8\ 15\ 6$ & $1\ 13\ 5$ & $+5$ & $1\ 11\ 58$ & $+2$ \\
11 & $10\ 34\ 3$ & $1\ 3\ 35$ & $+11$ & $1\ 25\ 51$ & $+10$ \\
12 & $11\ 18\ 32$ & $0\ 54\ 27$ & $+22$ & $1\ 39\ 3$ & $-1$ \\
13 & $8\ 54\ 0$ & $0\ 46\ 17$ & $+14$ & $1\ 50\ 41$ & $-9$ \\
14 & $8\ 24\ 44$ & $0\ 37\ 15$ & $+8$ & $2\ 3\ 27$ & $ .\quad . $ \\
15 & $10\ 54\ 28$ & $0\ 26\ 52$ & $+12$ & $2\ 17\ 57$ & $-8$ \\
17 & $10\ 23\ 9$ & $0\ 7\ 50$ & $+25$ & $2\ 44\ 12$ & $-20$ \\
18 & $8\ 38\ 17$ & $359\ 58\ 44$ & $-3$ & $2\ 56\ 35$ & $-16$ \\
21 & $8\ 30\ 54$ & $359\ 28\ 49$ & $-3$ & $3\ 36\ 50$ & $-4$ \\
21 & $10\ 9\ 32$ & $359\ 28\ 8$ & $-1$ & $3\ 37\ 45$ & $-1$ \\
23 & $13\ 25\ 57$ & $359\ 6\ 26$ & $+8$ & $4\ 6\ 32$ & $-5$ \\
24 & $8\ 27\ 37$ & $358\ 58\ 22$ & $+8$ & $4\ 17\ 12$ & $-10$ \\
25 & $8\ 41\ 38$ & $358\ 48\ 15$ & $+2$ & $4\ 30\ 32$ & $-22$ \\
\end{tabular}\endgroup
\end{center}

\Needspace{16\baselineskip}
\begin{center}
\textit{Braunschweiger Beobachtungen;} 
\end{center}

\begin{center}
\begingroup\fontsize{8}{10}\selectfont\setlength{\tabcolsep}{3pt}\renewcommand{\arraystretch}{1.12}
\begin{tabular}{r|r|r|r|r|r}
1804 & \shortstack{Mittlere Zeit in\\Braunschweig} & \shortstack{Berechnete gerade\\Aufsteigung} & \shortstack{Differenz\\R—B} & \shortstack{Berechnete\\Deklination} & \shortstack{Differenz\\R—B} \\ \hline
12. Sept. & $10^{u}\,35^{m}\,2^{s}$ & $0^\circ\ 54'\ 46''$ & $+15''$ & $1^\circ\ 38'\ 36''$ & $+21''$ \\
13 & $9\ 41\ 32$ & $0\ 46\ 2$ & $+38$ & $1\ 51\ 2$ & $+3$ \\
14 & $11\ 38\ 49$ & $0\ 35\ 58$ & $+21$ & $2\ 5\ 15$ & $+14$ \\
15 & $10\ 16\ 17$ & $0\ 27\ 10$ & $+17$ & $2\ 17\ 33$ & $-2$ \\
16 & $10\ 37\ 4$ & $0\ 17\ 28$ & $+11$ & $2\ 30\ 58$ & $-22$ \\
17 & $11\ 28\ 59$ & $0\ 7\ 25$ & $+2$ & $2\ 44\ 44$ & $+15$ \\
18 & $11\ 22\ 16$ & $359\ 57\ 38$ & $+12$ & $2\ 58\ 4$ & $+1$ \\
21 & $10\ 24\ 52$ & $359\ 28\ 5$ & $-1$ & $3\ 37\ 50$ & $-10$ \\
24 & $10\ 1\ 2$ & $358\ 57\ 45$ & $-8$ & $4\ 18\ 1$ & $-5$ \\
25 & $8\ 44\ 25$ & $358\ 48\ 6$ & $-6$ & $4\ 30\ 43$ & $-1$ \\
27 & $10\ 29\ 29$ & $358\ 27\ 8$ & $-12$ & $4\ 58\ 17$ & $+30$ \\
28 & $8\ 29\ 4$ & $358\ 17\ 50$ & $-30$ & $5\ 10\ 30$ & $+8$ \\
\end{tabular}\endgroup
\end{center}

Auch Dr. \textsc{Olbers} war so freundlich, uns seine weiteren Beobachtungen des merkwürdigen Fremdlings mitzuteilen:

\begin{center}
\begingroup\fontsize{9}{11}\selectfont\setlength{\tabcolsep}{3pt}\renewcommand{\arraystretch}{1.12}
\begin{tabular}{r|r|r|r}
1804 & \shortstack{Mittlere Zeit\\in\\Bremen} & \shortstack{Scheinbare\\gerade\\Aufsteigung} & \shortstack{Scheinbare\\südliche\\Deklination} \\ \hline
9. Sept. & $10^{u}\,48^{m}\,50^{s}$ & $1^\circ\ 20'\ 30''$ & $1^\circ\ 1'\ 5''$ \\
10 & $8\ 15\ 6$ & $1\ 12\ 55$ & $1\ 11\ 55$ \\
11 & $10\ 43\ 54$ & $1\ 3\ 20$ & $1\ 25\ 48$ \\
\end{tabular}\endgroup
\end{center}

Dr. \textsc{Olbers} schreibt dazu: „Diese Beobachtungen bedürfen noch einer kleinen Korrektur, da alle Vergleiche bis auf einen mit lichtschwachen Sternen aus der \textit{Hist. cél.} vorgenommen werden mussten. Da die Beobachtung von \textsc{Harding} und meine wahrscheinlich die ersten sind, die man an diesem neuen Planeten angestellt hat, wäre es sehr zu wünschen, dass Sie so freundlich wären, diese lichtschwachen Sterne von \textsc{La Lande} genauer zu bestimmen. Es sind folgende Sterne:

\Needspace{11\baselineskip}
\begin{center}\fontsize{9}{11}\selectfont
\begin{tabular}{l}
S. 119 Hist. cél.\\
Größe $8$\quad dritter Faden $0^u\ 8^m\ 53^s5$\\
\quad (soll sein $0^u\ 7^m\ 53^s5$)\quad Zenitdistanz $49^\circ\ 20'\ 6''$
\end{tabular}
\[
\text{S. 131 Hist. cél.}\left\{\begin{array}{lll}
8\ \text{—}&\text{mittlerer Faden}\ 0\ 3\ 40.5&\text{— —}\ 49\ 59\ 0\\
8\ \text{—}&\text{— —}\ 0\ 4\ 26.0&\text{— —}\ 50\ 15\ 12
\end{array}\right.
\]
\end{center}
\par
\SourceLeaf{258}%
\indent Auch ich habe diese drei Sterne so genau wie möglich bestimmt. Hier folgen ihre mittleren Positionen für den Anfang des Jahres 1804:

\begin{center}
\begingroup\fontsize{8}{10}\selectfont\setlength{\tabcolsep}{3pt}\renewcommand{\arraystretch}{1.12}
\begin{tabular}{r|r|r|r|r}
Größe & \shortstack{Mittlere gerade\\Aufsteigung\\1804} & \shortstack{Jährliche\\Änderung} & \shortstack{Mittlere südliche\\Deklination\\1804} & \shortstack{Jährliche\\Änderung} \\ \hline
$8$ & $1^\circ\ 0'\ 24''0$ & $+46''0$ & $1^\circ\ 7'\ 19''6$ & $-20''0$ \\
$8$ & $1\ 11\ 52.7$ & $+46.0$ & $1\ 23\ 30.9$ & $-20.0$ \\
$8$ & $2\ 0\ 20.9$ & $+46.0$ & $0\ 28\ 15.7$ & $-20.0$ \\
\end{tabular}\endgroup
\end{center}

In der Geschichte der Astronomie aller Zeiten und Nationen ist es beispiellos und zeugt von der glänzenden Epoche der heutigen Sternkunde in Deutschland, dass ein Planet vorhergesagt wurde und innerhalb von nur drei Wochen entdeckt und beobachtet, seine Bahn berechnet und sein künftiger Lauf vorgezeichnet wurde. All dies geschah jedoch durch die vereinten Kräfte von vier deutschen Astronomen, die das alles schon geleistet hatten, ehe die Nachricht von der Existenz dieses neuen Weltkörpers unsere eifersüchtigen Nachbarn erreicht hatte.

Zum Schluss folgt also eine Ephemeride des künftigen geozentrischen Laufs dieses neuen Planeten, die Dr. \textsc{Gauss} nach seinen obigen I. Elementen berechnet hat. So großer Verbesserungen diese Elemente auch noch bedürfen mögen, hofft Dr. \textsc{Gauss} doch zuversichtlich, dass sie diesen Lauf mehrere Wochen lang genau darstellen werden. Daher wird die folgende Ephemeride den weiter entfernten Astronomen sehr willkommen sein, die diesen neuen Himmelsgast noch nicht gefunden haben. Die Zeitpunkte sind für Mitternacht in Seeberg berechnet.

\begin{center}
\begingroup\fontsize{9}{11}\selectfont\setlength{\tabcolsep}{3pt}\renewcommand{\arraystretch}{1.12}
\begin{tabular}{r|r|r}
\shortstack{Mitternacht\\auf dem Seeberg} & \shortstack{Gerade Aufsteigung\\des neuen\\Planeten} & \shortstack{Südliche Deklination\\des neuen\\Planeten.} \\ \hline
5. Sept. & $1^\circ\ 52'$ & $0^\circ\ 12'$ \\
12 & $0\ 54$ & $1\ 39$ \\
19 & $359\ 47$ & $3\ 12$ \\
26 & $358\ 36$ & $4\ 55$ \\
30 & $357\ 56$ & $5\ 39$ \\
3. Okt. & $357\ 27$ & $6\ 17$ \\
6 & $357\ 0$ & $6\ 54$ \\
9 & $356\ 35$ & $7\ 29$ \\
12 & $356\ 12$ & $8\ 2$ \\
15 & $355\ 51$ & $8\ 33$ \\
\end{tabular}\endgroup
\end{center}
\begin{center}
\begingroup\fontsize{9}{11}\selectfont\setlength{\tabcolsep}{3pt}\renewcommand{\arraystretch}{1.12}
\begin{tabular}{r|r|r}
\shortstack{Mitternacht\\auf dem Seeberg} & \shortstack{Gerade Aufsteigung\\des neuen\\Planeten} & \shortstack{Südliche Deklination\\des neuen\\Planeten} \\ \hline
18. Okt. & $355^\circ\ 35'$ & $9^\circ\ 1'$ \\
21 & $355\ 21$ & $9\ 27$ \\
24 & $355\ 11$ & $9\ 50$ \\
27 & $355\ 4$ & $10\ 11$ \\
30 & $355\ 2$ & $10\ 28$ \\
2. Nov. & $355\ 4$ & $10\ 42$ \\
5 & $355\ 10$ & $10\ 54$ \\
8 & $355\ 20$ & $11\ 3$ \\
11 & $355\ 34$ & $11\ 8$ \\
14 & $355\ 53$ & $11\ 13$ \\
\end{tabular}\endgroup
\end{center}

\hfill[\textsc{von Zach.}]
